\newcommand{\matiere}{Chimie : les alcanes (correction)}
% Préambule commun aux fiches d'exercices et de correction
\documentclass[11pt,a4paper]{article}
\usepackage[utf8]{inputenc}
\usepackage[T1]{fontenc}
\usepackage{lmodern}
\usepackage[french]{babel}
\usepackage[a4paper,margin=2.2cm,headheight=14pt]{geometry}
\usepackage{amsmath,amssymb}
\usepackage{icomma}
\usepackage[version=4]{mhchem}
\usepackage{chemfig}
\usepackage{tikz}
\usetikzlibrary{decorations.pathmorphing,arrows.meta,calc,patterns,angles,quotes}
\usepackage{enumitem}
\setlist[itemize]{label=\textbullet,nosep}
\setlist[enumerate]{nosep}
\usepackage{array}
\usepackage[most]{tcolorbox}
\usepackage{fancyhdr}
\usepackage{lastpage}
\usepackage{needspace}
\newcommand{\vv}[1]{\overrightarrow{#1}}

\setlength{\parindent}{0pt}
\setlength{\parskip}{5pt}
\setchemfig{atom sep=2.1em}

\pagestyle{fancy}
\fancyhf{}
\fancyhead[L]{\small Préparation du devoir}
\fancyhead[R]{\small\matiere}
\fancyfoot[C]{\small Page \thepage\ sur \pageref{LastPage}}
\renewcommand{\headrulewidth}{0.4pt}

% Titre d'un exercice
\newcommand{\exercice}[1]{%
  \par\bigskip\needspace{9\baselineskip}
  \begin{tcolorbox}[colback=gray!12,colframe=black,boxrule=0.5pt,arc=1pt,
    left=6pt,right=6pt,top=2pt,bottom=2pt]
  \textbf{#1}
  \end{tcolorbox}\par\smallskip}

% Titre d'une question dans la correction
\newcommand{\question}[1]{\par\medskip\textbf{#1}\par}

% Résultat encadré dans la correction
\newtcolorbox{resultat}{colback=gray!6,colframe=black,boxrule=0.4pt,arc=1pt,
  left=6pt,right=6pt,top=2pt,bottom=2pt}

% Données encadrées
\newtcolorbox{donnees}{colback=white,colframe=black!60,boxrule=0.4pt,arc=1pt,
  left=6pt,right=6pt,top=2pt,bottom=2pt}

% Page de titre
\newcommand{\entete}[2]{%
  \begin{center}
  {\LARGE\bfseries #1}\\[4pt]
  {\large #2}\\[2pt]
  {\normalsize Classe de Première}
  \end{center}}

\begin{document}

\entete{Correction}{Chimie : les alcanes}

%-------------------------------------------------------------------------------
\exercice{Situation d'évaluation}

\question{1. Formule générale des alcanes}
La formule générale des alcanes est \ce{C_nH_{2n+2}}, où $n$ est le nombre d'atomes de carbone.

\question{2. Équation-bilan générale de la combustion complète}
\[
\ce{C_nH_{2n+2}} + \frac{3n+1}{2}\,\ce{O2} \longrightarrow n\,\ce{CO2} + (n+1)\,\ce{H2O}
\]

\question{3. Formule brute de l'alcane}
Quantités de matière des produits :
\[
n(\ce{CO2}) = \frac{m(\ce{CO2})}{M(\ce{CO2})} = \frac{3,52}{44} = 0,08 \text{ mol}
\qquad
n(\ce{H2O}) = \frac{m(\ce{H2O})}{M(\ce{H2O})} = \frac{1,8}{18} = 0,1 \text{ mol}
\]
D'après le bilan molaire de l'équation :
\[
\frac{n(\text{alcane})}{1} = \frac{n(\ce{CO2})}{n} = \frac{n(\ce{H2O})}{n+1}
\quad\Longrightarrow\quad
\frac{0,08}{n} = \frac{0,1}{n+1}
\]
\[
0,08\,(n+1) = 0,1\,n \;\Longrightarrow\; 0,08 = 0,02\,n \;\Longrightarrow\; n = 4.
\]
\begin{resultat}
L'alcane a pour formule brute \ce{C4H10} (butane). Comme sa chaîne est ramifiée, il s'agit du
2-méthylpropane : \chemfig{CH_3-CH(-[6]CH_3)-CH_3}
\end{resultat}

\question{4. Volume de dioxygène nécessaire}
Avec $n = 4$, l'équation-bilan s'écrit :
\[
\ce{C4H10 + \frac{13}{2} O2 -> 4 CO2 + 5 H2O}
\qquad\text{soit}\qquad
\ce{2 C4H10 + 13 O2 -> 8 CO2 + 10 H2O}
\]
Quantité de matière d'alcane brûlé ($M(\ce{C4H10}) = 4\times 12 + 10\times 1 = 58$ g/mol) :
\[
n(\ce{C4H10}) = \frac{1,16}{58} = 0,02 \text{ mol.}
\]
D'après le bilan molaire : $n(\ce{O2}) = \dfrac{13}{2}\, n(\ce{C4H10}) = 6,5 \times 0,02 = 0,13$ mol.
\[
V(\ce{O2}) = n(\ce{O2}) \times V_m = 0,13 \times 22,4 = 2,912 \text{ L.}
\]
\begin{resultat}
Le volume de dioxygène nécessaire est $V(\ce{O2}) \approx 2,9$ L.
\end{resultat}

%-------------------------------------------------------------------------------
\exercice{Exercice 1}

a. 2-chloro-3-méthylbutane :
\[\chemfig{CH_3-CH(-[6]Cl)-CH(-[6]CH_3)-CH_3}\]

b. 2,3-dichloro-2-méthylpentane :
\[\chemfig{CH_3-C(-[2]Cl)(-[6]CH_3)-CH(-[2]Cl)-CH_2-CH_3}\]

c. 2-bromo-4-chloro-3-éthyl-3-méthylhexane :
\[\chemfig{CH_3-CH(-[6]Br)-C(-[2]C_2H_5)(-[6]CH_3)-CH(-[6]Cl)-CH_2-CH_3}\]

d. 1,1,2,2-tétrafluoroéthane :
\[\chemfig{CH(-[3]F)(-[5]F)-CH(-[1]F)(-[7]F)}\]

%-------------------------------------------------------------------------------
\exercice{Exercice 2}

\renewcommand{\arraystretch}{1.3}
\begin{tabular}{|c|c|p{0.72\textwidth}|}
\hline
\textbf{N°} & \textbf{Réponse} & \textbf{Justification} \\ \hline
1) & \textbf{F} & Cette définition est celle d'un hydrocarbure. Un alcane est un hydrocarbure
  saturé : ses atomes sont liés uniquement par des liaisons simples. \\ \hline
2) & \textbf{V} & C'est la définition d'un alcane non cyclique. \\ \hline
3) & \textbf{V} & Un carbone lié à au moins trois autres carbones crée une ramification. \\ \hline
4) & \textbf{F} & Les cyclanes sont des alcanes dont la chaîne carbonée est fermée (cyclique). \\ \hline
5) & \textbf{F} & La formule générale des alcanes est \ce{C_nH_{2n+2}}. \\ \hline
6) & \textbf{V} & Toutes les liaisons sont simples : la molécule ne peut pas donner de réaction
  d'addition. \\ \hline
\end{tabular}

%-------------------------------------------------------------------------------
\exercice{Exercice 3}

\question{1. Équation-bilan}
\[
\ce{C3H8 + Cl2 ->[\text{lumière}] C3H7Cl + HCl}
\]

\question{2. Formules semi-développées et noms}
\begin{tabular}{@{}ll@{}}
$A$ : \chemfig{CH_3-CH_2-CH_2-Cl} & : 1-chloropropane \\[14pt]
$B$ : \chemfig{CH_3-CH(-[6]Cl)-CH_3} & : 2-chloropropane
\end{tabular}

%-------------------------------------------------------------------------------
\exercice{Exercice 4}

\question{1. Équation-bilan de la combustion complète des alcanes}
\[
\ce{C_nH_{2n+2}} + \frac{3n+1}{2}\,\ce{O2} \longrightarrow n\,\ce{CO2} + (n+1)\,\ce{H2O}
\]

\question{2. Formule brute de l'alcane $A$}
\[
n(\ce{CO2}) = \frac{m(\ce{CO2})}{M(\ce{CO2})} = \frac{11}{44} = 0,25 \text{ mol}
\qquad
n(\ce{H2O}) = \frac{m(\ce{H2O})}{M(\ce{H2O})} = \frac{5,4}{18} = 0,3 \text{ mol}
\]
D'après le bilan molaire :
\[
\frac{n(\ce{CO2})}{n} = \frac{n(\ce{H2O})}{n+1}
\;\Longrightarrow\; \frac{0,25}{n} = \frac{0,3}{n+1}
\;\Longrightarrow\; 0,25\,(n+1) = 0,3\,n
\;\Longrightarrow\; n = 5.
\]
\begin{resultat}
L'alcane $A$ a pour formule brute \ce{C5H12}.
\end{resultat}
\emph{Vérification :} $M(\ce{C5H12}) = 72$ g/mol, donc $n(A) = 3,6/72 = 0,05$ mol, et
$5 \times 0,05 = 0,25$ mol de \ce{CO2} : c'est bien cohérent.

\question{3. Formules semi-développées des isomères de \ce{C5H12}}
\begin{tabular}{@{}l@{\qquad}l@{}}
\chemfig{CH_3-CH_2-CH_2-CH_2-CH_3} & pentane \\[18pt]
\chemfig{CH_3-CH(-[6]CH_3)-CH_2-CH_3} & 2-méthylbutane \\[30pt]
\chemfig{CH_3-C(-[2]CH_3)(-[6]CH_3)-CH_3} & 2,2-diméthylpropane \\[20pt]
\end{tabular}

\question{4.1. Formule et nom de $A$}
La monochloration de $A$ ne donne qu'un seul produit : tous les atomes d'hydrogène de $A$ sont
équivalents. C'est le cas uniquement pour le 2,2-diméthylpropane, dont les quatre groupes \ce{CH3}
sont identiques.
\begin{resultat}
$A$ est le 2,2-diméthylpropane.
\end{resultat}

\question{4.2. Équation-bilan de la monochloration}
\[
\chemfig{CH_3-C(-[2]CH_3)(-[6]CH_3)-CH_3} \; + \; \ce{Cl2}
\;\xrightarrow{\text{lumière}}\;
\chemfig{CH_3-C(-[2]CH_3)(-[6]CH_3)-CH_2Cl} \; + \; \ce{HCl}
\]
Le produit obtenu est le 1-chloro-2,2-diméthylpropane.

%-------------------------------------------------------------------------------
\exercice{Exercice 5}
On note $V_1$ le volume initial de propane et $V_2$ le volume initial de butane (en cm$^3$).

\question{1. Équations-bilans de combustion}
\begin{tabular}{@{}lll@{}}
1.1. Propane : & \ce{C3H8 + 5 O2 -> 3 CO2 + 4 H2O} & volumes : $V_1$ ; $5V_1$ ; $3V_1$ \\[4pt]
1.2. Butane : & \ce{C4H10 + \frac{13}{2} O2 -> 4 CO2 + 5 H2O} & volumes : $V_2$ ; $\frac{13}{2}V_2$ ; $4V_2$
\end{tabular}

\question{2. Volumes initiaux de propane et de butane}
\begin{itemize}
  \item Volume du mélange : $V_1 + V_2 = 12$.
  \item Volume de \ce{CO2} formé : $3V_1 + 4V_2 = 42$.
  \item Volume de \ce{O2} consommé : $5V_1 + \frac{13}{2}V_2 = 100 - 31 = 69$.
\end{itemize}
De la première équation, $V_1 = 12 - V_2$. En remplaçant dans la deuxième :
\[
3(12 - V_2) + 4V_2 = 42 \;\Longrightarrow\; 36 + V_2 = 42 \;\Longrightarrow\; V_2 = 6 \text{ cm}^3,
\quad\text{puis}\quad V_1 = 12 - 6 = 6 \text{ cm}^3.
\]
\emph{Vérification avec le dioxygène :} $5 \times 6 + \frac{13}{2} \times 6 = 30 + 39 = 69$ cm$^3$.
\begin{resultat}
2.1. Volume initial du propane : $V_1 = 6$ cm$^3$. \qquad
2.2. Volume initial du butane : $V_2 = 6$ cm$^3$.
\end{resultat}

\question{3. Pourcentages en volume}
\[
\%\,\ce{C3H8} = \frac{V_1}{12} \times 100 = \frac{6}{12} \times 100 = 50\,\%
\qquad
\%\,\ce{C4H10} = \frac{V_2}{12} \times 100 = \frac{6}{12} \times 100 = 50\,\%
\]
\begin{resultat}
Le mélange initial contient $50\,\%$ de propane et $50\,\%$ de butane (en volume).
\end{resultat}

\end{document}
